% This template is used for all validation runs in the report (which can use different input datasets!)
\section{Validation run \detokenize{$<ContentVars['report_run_index']>$}}

Fig. \ref{fig:coverage\detokenize{$<ContentVars['report_run_index']>$}} shows the extent of validation run  \detokenize{$<ContentVars['report_run_index']>$}.
All statistics and plots in this chapter are based on data from the shown bounding box.

% A graphic created during the validation run data collection showing the
% covered area is included.
\begin{figure}[ht!]
    \centering
    \includegraphics[width=.6\textwidth]{"extent.png"}
    \caption{Area covered by all validation runs.}
    \label{fig:coverage\detokenize{$<ContentVars['report_run_index']>$}}
\end{figure}

\subsection{Validation Settings}

The following validation settings were used in the comparison against \detokenize{$<ContentVars['ConfigVars']['DS0']['name']>$}.
For a full list of used settings, see the validation run results page linked in Table \ref{tab:val_run_list}.
% Print a list of selected validation settings, information taken from the
% collected content variables (validation config section).
% This also shows some simple if-else logic in the f-string like placeholders.
\begin{itemize}
    \item Spatial reference: \textbf{\detokenize{$<ContentVars['ConfigVars']['spatial_reference']>$}}
    \item Temporal reference: \textbf{\detokenize{$<ContentVars['ConfigVars']['temporal_reference']>$}}
    \item Scaling reference: \textbf{\detokenize{$<ContentVars['ConfigVars']['scaling_reference']>$}}
    \item Temporal matching window was set to \textbf{\detokenize{$<ContentVars['ConfigVars']['temporal_matching']>$} hours}.
    \item Validation metrics \detokenize{$<'\\textbf{are}' if ContentVars['ConfigVars']['anomalies_method']!='none' else '\\textbf{are not}'>$} based on anomalies.
    \item Triple collocation analysis \detokenize{$<'\\textbf{was}' if ContentVars['ConfigVars']['metrics'][0]['value'] else '\\textbf{was not}'>$} performed.
\end{itemize}

\subsection{Results}

% Here we do some on-the-fly computation with numpy
\noindent SMOS L2 was compared against \detokenize{$<ContentVars['ConfigVars']['DS0']['name']>$} at \detokenize{$<ContentVars['NetcdfVars']['n_gpis']>$} locations.

\vspace{5mm}

\noindent Table \ref{tab:sumstats\detokenize{$<ContentVars['report_run_index']>$}} shows summary statistics in terms of number of
available observations / sample size for the metric computation ("\# Obs."), the correlation coefficient ("Pearson's R"),
Bias (difference in means) and unbiased RMSD. Additional metrics can be found in the validation results \href{\detokenize{$<ContentVars['NetcdfMetaVars']['url']>$}}{online}.

\begin{table}[htb!]
    \centering
    \caption{Summary statistics of SMOS L2 against \detokenize{$<ContentVars['ConfigVars']['DS0']['name']>$}.}
    \label{tab:sumstats\detokenize{$<ContentVars['report_run_index']>$}}
    \begin{tabular}{lrrrr}
        \toprule
        \textbf{Metric} & \textbf{\#Obs.} & \textbf{Pearson's R} & \textbf{Bias}  & \textbf{ubRMSD}  \\
        \midrule
        \textbf{Mean}   & \detokenize{$<ContentVars['SummaryVars']['MEAN_N. OBS.']>$}   & \detokenize{$<ContentVars['SummaryVars']['MEAN_PEARSON'S R']>$}     & \detokenize{$<ContentVars['SummaryVars']['MEAN_BIAS']>$}   & \detokenize{$<ContentVars['SummaryVars']['MEAN_UNBIASED ROOT-MEAN-SQUARE DEVIATION']>$}   \\
        \textbf{Median}   & \detokenize{$<ContentVars['SummaryVars']['MEDIAN_N. OBS.']>$}   & \detokenize{$<ContentVars['SummaryVars']['MEDIAN_PEARSON'S R']>$}     & \detokenize{$<ContentVars['SummaryVars']['MEDIAN_BIAS']>$}   & \detokenize{$<ContentVars['SummaryVars']['MEDIAN_UNBIASED ROOT-MEAN-SQUARE DEVIATION']>$}   \\
        \textbf{IQ-Range}   & \detokenize{$<ContentVars['SummaryVars']['IQ RANGE_N. OBS.']>$}   & \detokenize{$<ContentVars['SummaryVars']['IQ RANGE_PEARSON'S R']>$}     & \detokenize{$<ContentVars['SummaryVars']['IQ RANGE_BIAS']>$}   & \detokenize{$<ContentVars['SummaryVars']['IQ RANGE_UNBIASED ROOT-MEAN-SQUARE DEVIATION']>$}   \\
        \bottomrule
    \end{tabular}
\end{table}


\noindent For \detokenize{$<np.mean(ContentVars['NetcdfVars']['percent_ok']):.2f>$} \% of all tested locations, no validation metrics could be computed, because an error occurred
(compare Figure \ref{fig:status\detokenize{$<ContentVars['report_run_index']>$}}).

\begin{figure}[ht!]
\centering
\begin{minipage}[c]{0.3\textwidth}
    \centering
    \includegraphics[width=\textwidth]{./qa4sm_graphics/bulk_barplot_status.png}
    \subcaption{} % (a)
\end{minipage}
\begin{minipage}[c]{0.6\textwidth}
    \centering
    \includegraphics[width=\textwidth]{./qa4sm_graphics/bulk_overview_status.png}
    \subcaption{} % (b)
\end{minipage}
\caption{Number of successful test cases and errors raised (a), and their respective location (b) in the validation run against \detokenize{$<ContentVars['ConfigVars']['DS0']['name']>$}}
\label{fig:status\detokenize{$<ContentVars['report_run_index']>$}}
\end{figure}

\noindent Validation metrics are computed from common observation time stamps between SMOS L2 and \detokenize{$<ContentVars['ConfigVars']['DS0']['name']>$}.
A \detokenize{$<ContentVars['ConfigVars']['temporal_matching']>$}-hour temporal matching window is applied to each \detokenize{$<ContentVars['ConfigVars']['temporal_reference']>$} observation, to
find a corresponding data point in \detokenize{$<ContentVars['ConfigVars']['DS0']['name']>$}. The resulting sample size to compute all validation metrics is shown in Figure
\ref{fig:nobs\detokenize{$<ContentVars['report_run_index']>$}}.


\begin{figure}[ht!]
\centering
\begin{minipage}[c]{0.3\textwidth}
    \centering
    \includegraphics[width=\textwidth]{./qa4sm_graphics/bulk_boxplot_n_obs.png}
    \subcaption{} % (a)
\end{minipage}
\begin{minipage}[c]{0.6\textwidth}
    \centering
    \includegraphics[width=\textwidth]{./qa4sm_graphics/bulk_overview_n_obs.png}
    \subcaption{} % (b)
\end{minipage}
\caption{Validation sample size distribution (a), and by validation location (b) in the validation run against \detokenize{$<ContentVars['ConfigVars']['DS0']['name']>$}}
\label{fig:nobs\detokenize{$<ContentVars['report_run_index']>$}}
\end{figure}

\noindent Figure \ref{fig:Rbox\detokenize{$<ContentVars['report_run_index']>$}} show the correlation coefficients (Pearson's R) (a, b) and
p-values (c) between SMOS L2 and \detokenize{$<ContentVars['ConfigVars']['DS0']['name']>$} at all successfully tested locations.

\begin{figure}[ht!]
\centering
\begin{minipage}[c]{0.45\textwidth}
    \centering
    \includegraphics[width=\textwidth]{./qa4sm_graphics/bulk_boxplot_R.png}
    \subcaption{} % (a)
\end{minipage}%
\hfill
\begin{minipage}[c]{0.5\textwidth}
    \centering
    \includegraphics[width=\textwidth]{./qa4sm_graphics/bulk_overview_0-SMOS_L2_and_1-\detokenize{$<ContentVars['NetcdfMetaVars']['val_dc_dataset1']>$}_R.png}
    \subcaption{} % (b)
    \medskip
    \includegraphics[width=\textwidth]{./qa4sm_graphics/bulk_overview_0-SMOS_L2_and_1-\detokenize{$<ContentVars['NetcdfMetaVars']['val_dc_dataset1']>$}_p_R.png}
    \subcaption{} % (c)
\end{minipage}
\caption{Distribution of Pearson's R values across all locations (a),
         spatial overview of Pearson's R (b) and corresponding p-values (c)
         between common observations from SMOS L2 and
         \detokenize{$<ContentVars['ConfigVars']['DS0']['name']>$}}
\label{fig:Rbox\detokenize{$<ContentVars['report_run_index']>$}}
\end{figure}


